\clearpage
\SourceLeaf{161}\pdfbookmark[0]{Liber secundus.}{liber-secundus}
\begin{center}
{\Large LIBER SECUNDUS.}\par\bigskip
{\large INVESTIGATIO ORBITARUM CORPORUM COELESTIUM EX OBSERVATIONIBUS GEOCENTRICIS.}\par\bigskip
\rule{3em}{.3pt}\par\bigskip
\pdfbookmark[1]{Sectio prima.}{liber-secundus-sectio-prima}
{\large SECTIO PRIMA.}\par\medskip
\emph{Determinatio orbitae e tribus observationibus completis.}\par\bigskip
\hypertarget{native-161-article115}{}\pdfbookmark[2]{115.}{article-115}115.
\end{center}
Ad determinationem completam motus corporis coelestis in orbita sua requiruntur elementa \emph{septem}, quorum autem numerus uno minor evadit, si corporis massa vel cognita est vel negligitur; haec licentia vix evitari poterit in determinatione orbitae penitus adhuc incognitae, ubi omnes quantitates ordinis perturbationum tantisper seponere oportet, donec massae a quibus pendent aliunde innotuerint. Quamobrem in disquisitione praesente massa corporis neglecta elementorum numerum ad sex reducimus, patetque adeo, ad determinationem orbitae incognitae totidem quantitates ab elementis pendentes ab invicem vero independentes requiri. Quae quantitates nequeunt esse nisi loca corporis coelestis e terra observata, quae singula quum bina data subministrent, puta longitudinem et latitudinem, vel ascensionem rectam et declinationem, simplicissimum utique erit, \emph{tria loca geocentrica} adoptare, quae generaliter loquendo sex elementis incognitis determinandis sufficient. Hoc problema \SourceLeaf{162}tamquam gravissimum huius operis spectandum erit, summaque ideo cura in hac Sectione pertractabitur.

Verum enim vero in casu speciali, ubi planum orbitae cum ecliptica coincidit adeoque omnes latitudines tum heliocentricae tum geocentricae natura sua evanescunt, tres latitudines geocentricas evanescentes haud amplius considerare licet tamquam tria data ab invicem independentia: tunc igitur problema istud indeterminatum maneret, tribusque locis geocentricis per orbitas infinite multas satisfieri posset. In tali itaque casu necessario quatuor longitudines geocentricas datas esse oportet, ut quatuor elementa incognita reliqua (excidentibus inclinatione orbitae et longitudine nodi) determinare liceat. Etiamsi vero per principium indiscernibilium haud exspectandum sit, talem casum in rerum natura unquam se oblaturum esse, tamen facile praesumitur, problema, quod in orbita cum plano eclipticae omnino coincidente absolute indeterminatum fit, \emph{in orbitis perparum ad eclipticam inclinatis} propter observationum praecisionem limitatam tantum non indeterminatum manere debere, ubi vel levissimi observationum errores incognitarum determinationem penitus turbare valent. Quamobrem ut huic quoque casui consulamus, alia sex data eligere oportebit: ad quem finem in Sectione secunda orbitam incognitam e quatuor observationibus determinare docebimus, quarum duae quidem completae sint, duae reliquae autem incompletae, latitudinibus vel declinationibus deficientibus.

Denique quum omnes observationes nostrae propter instrumentorum sensuumque imperfectionem non sint nisi approximationes ad veritatem, orbita, sex tantum datis absolute necessariis superstructa, erroribus considerabilibus adhuc obnoxia esse poterit. Quos ut quantum quidem licet extenuemus, summamque adeo praecisionem possibilem attingamus, via alia non dabitur, nisi ut observationes perfectissimas quam plurimas congeramus, elementaque ita perpoliamus, ut non quidem his vel illis praecisione absoluta satisfaciant, sed cum cunctis quam optime conspirent. Quonam pacto talem consensum, si nullibi absolutum tamen ubique quam arctissimum, secundum principia calculi probabilitatis obtinere liceat, in Sectione tertia ostendemus.

Hoc itaque modo determinatio orbitarum, quatenus corpora coelestia secundum leges \textsc{Kepler}i in ipsis moventur, ad omnem quae desiderari potest perfectionem evecta erit. Ultimam quidem expolitionem tunc demum suscipere \SourceLeaf{163}licebit, ubi etiam perturbationes, quas planetae reliqui motui inducunt, ad calculum erunt revocatae: quarum rationem quomodo habere oporteat, quantum quidem ad institutum nostrum pertinere videbitur, in Sectione quarta breviter indicabimus.
